\chapter{Representations and Characters}

\begin{center}\itshape
The group $\mathfrak S_4$ has $24$ elements, and $24=1+1+4+9+9$.\\
The group $A_5$ has $60$ elements, and $60=1+9+9+16+25$.\\
Why should the order of a group be a sum of squares?
\end{center}

The symmetric group on three letters has six elements. Yet
\[
6=1^2+1^2+2^2.
\]
Those squares are not chosen to make the arithmetic work. They are the dimensions of three irreducible complex representations of $\mathfrak S_3$. For the cube's rotation group $\mathfrak S_4$, the same rule becomes
\[
24=1^2+1^2+2^2+3^2+3^2.
\]

A \emph{representation} lets a group act by invertible matrices. Matrices are easier to compute with than an abstract multiplication law: we can add them, take traces, and use linear algebra. A \emph{character} records the trace of each matrix. At first, throwing away almost every entry seems like a terrible loss of information. The surprise is that characters retain enough information to classify irreducible complex representations of a finite group.

The reason is an averaging trick. Average a linear map over a finite group and it begins to respect the group action. That small move leads to Maschke's theorem, Schur's lemma, and the orthogonality relations for characters.

In the abelian case the picture becomes especially simple. Every irreducible complex representation is one-dimensional, so a character is just a homomorphism into the unit circle. The corresponding change of basis is the finite Fourier transform. You have already seen its roots of unity in earlier volumes; now they appear as the coordinates naturally chosen by symmetry.

We will compute tables for $\mathfrak S_4$ and $A_5$, then see how the abelian case reconnects with modular arithmetic and, eventually, the analysis of random walks.

\section{Representations}

\begin{definition}
A \emph{representation} of $G$ is a homomorphism $\rho:G\to GL(V)$, where $V$ is a complex vector space of dimension $d\ge1$, the \emph{dimension} (or degree) of $\rho$; we also say that $V$ is a representation, and write $gv$ for $\rho(g)v$. A subspace $W$ is \emph{stable} if $gW\subseteq W$ for all $g$; it is then a \emph{subrepresentation}. $V$ is \emph{irreducible} if its only stable subspaces are $0$ and $V$. A linear map $f:V\to V'$ between representations is an \emph{intertwiner} if $f(gv)=gf(v)$ for all $g,v$; an invertible one is an \emph{isomorphism}.
\end{definition}

\begin{example}\label{ex:reps}
\begin{enumerate}[label=(\arabic*)]
\item The \emph{trivial} representation: $V=\C$, $\rho(g)=1$. The \emph{sign} representation of $\mathfrak S_n$: $\rho(\sigma)=\varepsilon(\sigma)\in\{\pm1\}=GL_1(\C)\cap\R$.
\item Every one-dimensional representation is a homomorphism $\chi:G\to\C^*$, and is irreducible. For $G=\Z/n\Z$, $\chi_j(k)=e^{2i\pi jk/n}$, $0\le j<n$.
\item The \emph{permutation representation} of $\mathfrak S_n$ on $\C^n$: $\sigma e_i=e_{\sigma(i)}$. The line $\C(e_1+\dots+e_n)$ and the hyperplane $\{\sum x_i=0\}$ are stable.
\item More generally, if $G$ acts on a finite set $X$ (Chapter~3), $G$ acts on the space $\C^X$ with basis $(e_x)_{x\in X}$ by $ge_x=e_{gx}$. With $X=G$ acting on itself by left multiplication, this is the \emph{regular representation} $\C[G]$, of dimension $|G|$.
\item The symmetries of an equilateral triangle centred at $0$ give a homomorphism $\mathfrak S_3\cong D_3\to GL_2(\R)\subseteq GL_2(\C)$. It is irreducible: a stable line would be a common eigenline of the rotation by $120^\circ$, which has no real eigenvector, and of a reflection; over $\C$ the rotation's eigenlines are spanned by $(1,\mp i)$, and the reflection in the $x$-axis swaps them.
\end{enumerate}
\end{example}

\begin{think}
$\mathfrak S_3$ has order $6$, and the trivial and sign representations have dimension $1$. If $6$ is a sum of the squares of the dimensions of all irreducible representations, what is missing? Which representation above has that dimension?
\end{think}

\begin{theorem}[Maschke]\label{thm:maschke}
Every representation $V$ of $G$ has a $G$-invariant Hermitian inner product: $\langle gv,gw\rangle=\langle v,w\rangle$. Consequently every stable subspace $W$ has a stable complement $W^\perp$, and $V$ is a direct sum of irreducible subrepresentations.
\end{theorem}
\begin{proof}
Start with an arbitrary Hermitian inner product $(v,w)$ on the complex vector space $V$. It may not be preserved by the representation. We make an invariant one by averaging over the group:
\[
\langle v,w\rangle=\frac1{|G|}\sum_{h\in G}(hv,hw).
\]
This is again Hermitian. It is positive definite because, for $v\ne0$, every term $(hv,hv)$ is positive: $h$ acts invertibly, so $hv\ne0$. Thus it is an inner product.

For any $g\in G$, replace $h$ by $hg$ in the sum. Right multiplication permutes the elements of $G$, so
\[
\langle gv,gw\rangle
=\frac1{|G|}\sum_{h\in G}(hgv,hgw)
=\langle v,w\rangle.
\]
The new inner product is invariant under $G$.

Let $W\subseteq V$ be an invariant subspace. We claim its orthogonal complement is invariant too. Take $v\in W^\perp$ and $w\in W$. Using invariance of the inner product,
\[
\langle gv,w\rangle=\langle v,g^{-1}w\rangle.
\]
Because $W$ is invariant, $g^{-1}w\in W$, so the right side is $0$. This holds for every $w\in W$, hence $gv\in W^\perp$.

We have obtained a direct sum $V=W\oplus W^\perp$ of invariant subspaces. If either is reducible, repeat the argument inside it. Their dimensions decrease, so the process ends with a direct sum of irreducible representations.
\end{proof}

\paragraph{Why this works.}
The obstacle in decomposing a representation is that an arbitrary linear complement to an invariant subspace need not be invariant. Start with any projection onto the invariant subspace and average its conjugates by the finite group. The average is still a projection onto the same subspace, because every conjugate acts as the identity there, and now it commutes with every group element. Its kernel is consequently invariant. The division by $|G|$ is the essential hypothesis: it works over $\C$ but can fail in a field whose characteristic divides $|G|$.

So each $\rho(g)$ is unitary for $\langle\cdot,\cdot\rangle$. Since $g^n=e$ for $n=|G|$, $\rho(g)$ is a unitary matrix with $\rho(g)^n=I$: it is diagonalisable in an orthonormal basis (spectral theorem, Volume~1), with eigenvalues that are $n$-th roots of unity.

\begin{lemma}[Schur]\label{lem:schur}
Let $V,V'$ be irreducible and $f:V\to V'$ an intertwiner. Then $f=0$ or $f$ is an isomorphism. If $V=V'$, then $f=\lambda\,\mathrm{id}$ for some $\lambda\in\C$.
\end{lemma}
\begin{proof}
$\ker f$ is stable ($f(v)=0\Rightarrow f(gv)=gf(v)=0$), so it is $0$ or $V$; $f(V)$ is stable, so it is $0$ or $V'$. If $f\ne0$, then $\ker f=0$ and $f(V)=V'$. If $V=V'$, let $\lambda$ be an eigenvalue of $f$ (we are over $\C$). Then $f-\lambda\,\mathrm{id}$ is an intertwiner that is not injective, hence $0$.
\end{proof}

\begin{corollary}\label{cor:abelianirr}
If $G$ is abelian, every irreducible representation has dimension $1$.
\end{corollary}
\begin{proof}
For $h\in G$, $\rho(h)$ commutes with all $\rho(g)$, so it is an intertwiner $V\to V$, hence a scalar. Then every line is stable, and $V$, irreducible, is a line.
\end{proof}

A representation stores a matrix for each group element. Taking traces reduces each matrix to one number, and those numbers stay constant on conjugacy classes. Orthogonality will show why this surprisingly short record can distinguish irreducible representations.

\section{Characters and orthogonality}

\begin{definition}
The \emph{character} of a representation $\rho$ is $\chi_\rho(g)=\operatorname{tr}\rho(g)$. A \emph{class function} is a function $G\to\C$ constant on conjugacy classes; they form a vector space of dimension the number of classes, with the inner product
\[
\langle\varphi,\psi\rangle=\frac1{|G|}\sum_{g\in G}\varphi(g)\overline{\psi(g)}.
\]
\end{definition}

\begin{proposition}\label{prop:charbasic}
\begin{enumerate}[label=(\alph*)]
\item $\chi(e)=\dim V$; $\chi$ is a class function; isomorphic representations have the same character.
\item $\chi(g^{-1})=\overline{\chi(g)}$, and $|\chi(g)|\le\chi(e)$.
\item $\chi_{V\oplus W}=\chi_V+\chi_W$.
\item For a permutation representation $\C^X$, $\chi(g)$ is the number of fixed points of $g$ on $X$. In particular the character of $\C[G]$ is $|G|$ at $e$ and $0$ elsewhere.
\end{enumerate}
\end{proposition}
\begin{proof}
(a) $\operatorname{tr}I=\dim V$, and $\operatorname{tr}(ABA^{-1})=\operatorname{tr}B$ applies both to $\rho(hgh^{-1})=\rho(h)\rho(g)\rho(h)^{-1}$ and to $\rho'(g)=f\rho(g)f^{-1}$. (b) The eigenvalues $\lambda_j$ of $\rho(g)$ are roots of unity, so $\lambda_j^{-1}=\overline{\lambda_j}$; those of $\rho(g^{-1})$ are the $\lambda_j^{-1}$; and $|\sum\lambda_j|\le\sum|\lambda_j|=\dim V$. (c) In a basis adapted to $V\oplus W$ the matrices are block diagonal. (d) $\rho(g)$ is a permutation matrix, whose diagonal entry at $e_x$ is $1$ iff $gx=x$. Left multiplication by $g\ne e$ has no fixed point.
\end{proof}

\begin{theorem}[orthogonality of characters]\label{thm:orth}
Let $\chi,\chi'$ be the characters of irreducible representations $V,V'$. Then $\langle\chi,\chi'\rangle=1$ if $V\cong V'$, and $0$ otherwise.
\end{theorem}
\begin{proof}
For any linear map $u:V\to V'$, the average
\[
\tilde u=\frac1{|G|}\sum_{g\in G}\rho'(g)\,u\,\rho(g)^{-1}
\]
is an intertwiner:
\[
\rho'(h)\,\tilde u\,\rho(h)^{-1}=\frac1{|G|}\sum_g\rho'(hg)\,u\,\rho(hg)^{-1}=\tilde u.
\]
 Choose bases, write matrices with entries $\rho(g)_{ij}$, and take for $u$ the matrix unit $E_{jk}$ (entry $1$ at $(j,k)$). The $(i,l)$ entry of $\tilde u$ is $\frac1{|G|}\sum_g\rho'(g)_{ij}\rho(g^{-1})_{kl}$.

If $V\not\cong V'$, Schur gives $\tilde u=0$ for every $u$, so all these sums vanish; put $i=j$, $k=l$ and sum over $j$ and $k$: $\frac1{|G|}\sum_g\chi'(g)\chi(g^{-1})=0$, i.e.\ $\langle\chi',\chi\rangle=0$ by Proposition~\ref{prop:charbasic}(b).

If $V\cong V'$, the characters are equal, so we may take $V'=V$, $\rho'=\rho$. Schur gives $\tilde u=\lambda I$, and taking traces, $\operatorname{tr}\tilde u=\operatorname{tr}u$, so $\lambda=\operatorname{tr}(u)/d$. For $u=E_{jk}$, $\operatorname{tr}u=\delta_{jk}$, hence $\frac1{|G|}\sum_g\rho(g)_{ij}\rho(g^{-1})_{kl}=\delta_{il}\delta_{jk}/d$. Put $i=j$, $k=l$ and sum: $\frac1{|G|}\sum_g\chi(g)\chi(g^{-1})=\sum_{j,k}\delta_{jk}/d=d/d=1$.
\end{proof}

\paragraph{Why this works.}
Trace is invariant under a change of basis, but that alone does not explain the strong orthogonality formula. The averaging operation takes an arbitrary linear map $V\to V'$ and turns it into one that respects the group action. Schur's lemma then says that the average is zero between different irreducible representations, and a scalar multiple of the identity when they are the same. Summing the diagonal matrix coefficients converts that statement into an inner product of character values. This is why a short list of traces can recognise a whole irreducible representation.

\begin{corollary}\label{cor:mult}
Let $V_1,\dots,V_r$ be pairwise non-isomorphic irreducible representations with characters $\chi_1,\dots,\chi_r$, and $V$ any representation.
\begin{enumerate}[label=(\alph*)]
\item In any decomposition of $V$ into irreducibles, the number of summands isomorphic to $V_i$ is $\langle\chi_V,\chi_i\rangle$. So two representations with the same character are isomorphic.
\item $V$ is irreducible iff $\langle\chi_V,\chi_V\rangle=1$.
\item (Regular representation) Every irreducible $V_i$ occurs in $\C[G]$, $\dim V_i$ times. Hence there are finitely many irreducible representations up to isomorphism, and if $V_1,\dots,V_r$ is the full list, $\sum_i(\dim V_i)^2=|G|$.
\end{enumerate}
\end{corollary}
\begin{proof}
(a) If $V\cong\bigoplus W_j$ with irreducible $W_j$, then $\chi_V=\sum\chi_{W_j}$ and, by Theorem~\ref{thm:orth}, $\langle\chi_V,\chi_i\rangle$ counts the $j$ with $W_j\cong V_i$. A character thus determines the multiplicities, hence the representation up to isomorphism (Maschke).

(b) If $V\cong\bigoplus V_i^{m_i}$ (all irreducibles that occur), $\langle\chi_V,\chi_V\rangle=\sum m_i^2$, which is $1$ iff $V$ is a single irreducible.

(c) By Proposition~\ref{prop:charbasic}(d), $\langle\chi_{\C[G]},\chi_i\rangle=\frac1{|G|}|G|\,\overline{\chi_i(e)}=\dim V_i$. Every irreducible occurs in the finite-dimensional $\C[G]$, so there are finitely many, and comparing dimensions in $\C[G]\cong\bigoplus V_i^{\dim V_i}$ gives $|G|=\sum(\dim V_i)^2$.
\end{proof}

\begin{theorem}\label{thm:classes}
The irreducible characters form an orthonormal basis of the space of class functions. In particular, the number of irreducible representations of $G$ is the number of its conjugacy classes.
\end{theorem}
\begin{proof}
They are orthonormal (Theorem~\ref{thm:orth}), hence linearly independent; it remains to show that a class function $f$ with $\langle f,\chi_i\rangle=0$ for all $i$ is $0$. For a representation $\rho$, let $T_\rho=\sum_g\overline{f(g)}\rho(g)$. As $\overline f$ is constant on classes, $\rho(h)T_\rho\rho(h)^{-1}=\sum_g\overline{f(g)}\rho(hgh^{-1})=T_\rho$, so $T_\rho$ is an intertwiner. On an irreducible $V_i$ it is a scalar $\lambda$ (Schur), with $\lambda\dim V_i=\operatorname{tr}T_\rho=\sum_g\overline{f(g)}\chi_i(g)=|G|\,\overline{\langle f,\chi_i\rangle}=0$. So $T_\rho=0$ on every irreducible, hence on every representation (Maschke), in particular on $\C[G]$. But $T_{\C[G]}e_e=\sum_g\overline{f(g)}e_g$, so all $f(g)=0$.
\end{proof}

\begin{corollary}[column orthogonality]\label{cor:columns}
Let $\chi_1,\dots,\chi_r$ be the irreducible characters, and $g,h\in G$. Then $\sum_i\chi_i(g)\overline{\chi_i(h)}$ is $|C_G(g)|$ if $g$ and $h$ are conjugate, and $0$ otherwise.
\end{corollary}
\begin{proof}
Let $C_1,\dots,C_r$ be the classes (there are $r$ of them, Theorem~\ref{thm:classes}), and $M$ the $r\times r$ matrix $M_{ij}=\sqrt{|C_j|/|G|}\,\chi_i(C_j)$. Row orthogonality says $M\overline M^{\,T}=I$. A square matrix with a right inverse is invertible with the same inverse, so $\overline M^{\,T}M=I$: $\sum_i\frac{\sqrt{|C_j||C_k|}}{|G|}\overline{\chi_i(C_j)}\chi_i(C_k)=\delta_{jk}$. With $|G|/|C_j|=|C_G(g)|$ for $g\in C_j$ (orbit--stabiliser, Theorem~\ref{thm:orbstab}), this is the claim.
\end{proof}

The irreducible characters form an orthonormal basis for class functions. Putting their values on conjugacy classes into a table makes the orthogonality relations visible and lets us determine unknown entries from a few constraints.

\section{Character tables}

The \emph{character table} of $G$ lists the values of the irreducible characters on the conjugacy classes; it is square by Theorem~\ref{thm:classes}. Three tools fill it in: one-dimensional representations, permutation representations (whose characters count fixed points), and the orthogonality relations. Two more constructions help. If $\chi$ is a character and $\lambda$ a one-dimensional character, $\lambda\chi$ is a character (of $\rho'(g)=\lambda(g)\rho(g)$ on the same space), irreducible with $\chi$ since $\langle\lambda\chi,\lambda\chi\rangle=\langle\chi,\chi\rangle$ ($|\lambda(g)|=1$). And if $N\trianglelefteq G$, every representation of $G/N$ gives one of $G$ by composing with $G\to G/N$; it is irreducible iff the original is, since the two have the same matrices.

\begin{example}[$\mathfrak S_3$]\label{ex:s3}
Three classes (identity, $3$ transpositions, $2$ three-cycles), so three irreducibles, with $\sum d_i^2=6$: dimensions $1,1,2$. The permutation character of $\C^3$ is $(3,1,0)$ (fixed points); minus the trivial character it gives $\chi_2=(2,0,-1)$, and $\langle\chi_2,\chi_2\rangle=\frac{4+0+2\cdot1}6=1$, so it is irreducible: it is the triangle representation of Example~\ref{ex:reps}(5).
\end{example}

\begin{example}[$\mathfrak S_4$]\label{ex:s4}
Five classes, of sizes $1,6,3,8,6$ for $e$, $(12)$, $(12)(34)$, $(123)$, $(1234)$ (Chapter~4). The permutation character of $\C^4$ is $(4,2,0,1,0)$; minus the trivial one, $\chi_3=(3,1,-1,0,-1)$, of norm $\frac{9+6+3+0+6}{24}=1$. Then $\chi_3'=\varepsilon\chi_3=(3,-1,-1,0,1)$ is another irreducible character. The quotient $\mathfrak S_4/V_4\cong\mathfrak S_3$ (Chapter~4) turns $\chi_2$ of $\mathfrak S_3$ into $(2,0,2,-1,0)$: the double transpositions lie in $V_4$, and the image of a $4$-cycle in $\mathfrak S_3$ is a transposition, since its square lies in $V_4$ but it does not. With the trivial and sign characters, these are five distinct irreducible characters, and $1+1+4+9+9=24$: the table is complete (Figure~\ref{fig:tables}). The two $3$-dimensional ones are the rotations of the cube ($\chi_3'$, through the isomorphism of Chapter~1) and the symmetries of the regular tetrahedron with vertices $e_1,\dots,e_4$ ($\chi_3$).
\end{example}

\begin{figure}[ht]
\centering
\fig{ch11-tables}
\caption{The character tables of $\mathfrak S_3$, $\mathfrak S_4$ and $A_5$. Columns are the conjugacy classes (a representative, and the size of the class in grey), rows the irreducible characters; each cell is coloured by $\chi(g)/\chi(e)$, from charcoal ($-1$) through white ($0$) to coral ($1$). In the table of $A_5$, $\varphi=\frac{1+\sqrt5}2$ is the golden ratio. Rows are orthonormal for the weighted inner product, and columns are orthogonal (Corollary~\ref{cor:columns}).}
\label{fig:tables}
\end{figure}

\begin{remark}
Characters prove theorems about groups that have nothing to do with linear algebra. Burnside (1904) showed with them that a group of order $p^aq^b$ ($p,q$ primes) is solvable, and Frobenius that in a group acting transitively on a set where no non-identity element fixes two points, the identity together with the elements fixing no point form a normal subgroup. No proof of the latter avoiding characters is known. The dimensions of the irreducible representations always divide $|G|$ (Exercise~\ref{ex:divides}).
\end{remark}

For an abelian group, every conjugacy class has one element and every irreducible complex representation has dimension one. Its character table is therefore the Fourier matrix. The abstract theory becomes the familiar arithmetic of roots of unity.

\section{Abelian groups: Fourier analysis}\label{sec:fourierab}

Let $A$ be a finite abelian group. By Corollary~\ref{cor:abelianirr} and Theorem~\ref{thm:classes} (each element is its own class), $A$ has exactly $|A|$ irreducible characters, all one-dimensional: the homomorphisms $\chi:A\to\C^*$. They form a group $\widehat A$ under pointwise multiplication (a product of homomorphisms to the abelian group $\C^*$ is one), the \emph{dual group}, and $|\widehat A|=|A|$. For $A=\Z/n\Z$, a homomorphism is determined by $\chi(1)$, an $n$-th root of unity, so $\widehat A=\{\chi_j:k\mapsto e^{2i\pi jk/n}\}\cong\Z/n\Z$.

\begin{theorem}[Fourier analysis on a finite abelian group]\label{thm:fourier}
For $f:A\to\C$ define $\hat f(\chi)=\langle f,\chi\rangle=\frac1{|A|}\sum_{a}f(a)\overline{\chi(a)}$. Then
\[
\begin{gathered}f=\sum_{\chi\in\widehat A}\hat f(\chi)\,\chi,\qq\frac1{|A|}\sum_{a\in A}|f(a)|^2=\sum_{\chi\in\widehat A}|\hat f(\chi)|^2,\end{gathered}
\]
and for $a\in A$,
\[
\frac1{|A|}\sum_{\chi\in\widehat A}\chi(a)=\begin{cases}1&\text{if }a=0,\\0&\text{otherwise.}\end{cases}
\]
\end{theorem}
\begin{proof}
All functions are class functions, and the characters form an orthonormal basis of them (Theorem~\ref{thm:classes}); the first formula is the expansion in an orthonormal basis, the second is Pythagoras (Parseval). The third is Corollary~\ref{cor:columns} with $h=0$: the centraliser of $a$ is $A$, and $a$ is conjugate to $0$ only if $a=0$.
\end{proof}

For $A=\Z/n\Z$, $\hat f(\chi_j)=\frac1n\sum_kf(k)e^{-2i\pi jk/n}$ is the \emph{discrete Fourier transform}. It turns convolution into multiplication (Exercise~\ref{ex:circulant}), and the fast Fourier transform of Cooley and Tukey computes it in $O(n\log n)$ operations instead of $O(n^2)$, by splitting $\Z/n\Z$ along a subgroup when $n=2^k$.

\subsection*{Dirichlet characters}

Let $q\ge1$. A \emph{Dirichlet character} modulo $q$ is a function $\chi:\Z\to\C$ obtained from a character $\chi$ of $(\Z/q\Z)^*$ by $\chi(n)=\chi(n\bmod q)$ if $\gcd(n,q)=1$, and $\chi(n)=0$ otherwise. It is periodic of period $q$ and \emph{completely multiplicative}: $\chi(mn)=\chi(m)\chi(n)$ for all $m,n$ (if $m$ or $n$ is not prime to $q$, both sides are $0$). There are $\varphi(q)$ of them; the \emph{principal} character $\chi_0$ is $1$ on integers prime to $q$.

\begin{proposition}\label{prop:dirichletorth}
Let $a$ be prime to $q$. For every integer $n$,
\[
\frac1{\varphi(q)}\sum_{\chi\bmod q}\overline{\chi(a)}\,\chi(n)=\begin{cases}1&\text{if }n\equiv a\ (q),\\0&\text{otherwise.}\end{cases}
\]
\end{proposition}
\begin{proof}
If $\gcd(n,q)>1$, every term is $0$, and $n\not\equiv a$. Otherwise $\overline{\chi(a)}\chi(n)=\chi(a^{-1}n)$ in $A=(\Z/q\Z)^*$ ($|\chi(a)|=1$, so $\overline{\chi(a)}=\chi(a)^{-1}=\chi(a^{-1})$), and Theorem~\ref{thm:fourier}, written multiplicatively, gives $1$ if $a^{-1}n=1$ in $A$ and $0$ otherwise.
\end{proof}

This identity is a \emph{filter}: to study the primes $p\equiv a\pmod q$, sum $\frac1{\varphi(q)}\overline{\chi(a)}\chi(p)$ over all primes and all $\chi$, then treat each $\chi$ separately, which is possible because $\chi$ is multiplicative and has an Euler product. This is the plan of Dirichlet's theorem, carried out in Volume~6. For $q=p$ an odd prime, $(\Z/p\Z)^*$ is cyclic of even order, so it has exactly one character of order $2$, which must be the Legendre symbol $\bigl(\frac np\bigr)$ of Volume~2: Euler's criterion $\bigl(\frac np\bigr)\equiv n^{(p-1)/2}$ shows that it is multiplicative.

\begin{example}\label{ex:dir5}
Modulo $5$, $2$ generates $(\Z/5\Z)^*=\{1,2,4,3\}=\{2^0,2^1,2^2,2^3\}$, and the four characters are determined by $\chi(2)\in\{1,i,-1,-i\}$:
\begin{center}
\small
\begin{tabular}{c|cccc}
$n\bmod5$ & $1$ & $2$ & $3$ & $4$\\\hline
$\chi_0$ & $1$ & $1$ & $1$ & $1$\\
$\chi_1$ & $1$ & $i$ & $-i$ & $-1$\\
$\chi_2$ & $1$ & $-1$ & $-1$ & $1$\\
$\chi_3$ & $1$ & $-i$ & $i$ & $-1$
\end{tabular}
\end{center}
$\chi_2$ is the Legendre symbol $\bigl(\frac n5\bigr)$: the squares modulo $5$ are $1$ and $4$.
\end{example}

\begin{beauty}
The same few lines of linear algebra describe the vibrations of a symmetric molecule (chemists read the character table of its symmetry group), the frequencies of a sampled signal (the Fourier transform on $\Z/n\Z$), the mixing time of card shuffles (Diaconis and Shahshahani used the characters of $\mathfrak S_n$ to show that about $\frac12n\log n$ random transpositions mix a deck of $n$ cards, Volume~4), and the distribution of primes in progressions (Dirichlet characters, Volume~6). Fourier analysis on $\R$ and on the circle, met in Volume~5, is the same theory for infinite abelian groups; for non-abelian Lie groups such as $SO(3)$, the next chapter's rotations, it becomes the theory of spherical harmonics and of spin in quantum mechanics.
\end{beauty}

\section{Exercises}

\begin{exercise}[$\star$]\label{ex:pairs}
$\mathfrak S_4$ acts on the $6$ two-element subsets of $\{1,2,3,4\}$. Decompose this permutation representation into irreducibles.
\end{exercise}

\begin{exercise}[$\star$]\label{ex:onedim}
Show that the one-dimensional representations of $G$ correspond to the homomorphisms $G/G'\to\C^*$, where $G'$ is the derived subgroup (Exercise~\ref{ex:derived}), so that there are $[G:G']$ of them. How many do $\mathfrak S_n$ ($n\ge2$), $A_4$ and $A_5$ have?
\end{exercise}

\begin{exercise}[$\star\star$]\label{ex:a4}
Find the character table of $A_4$. (Its classes have sizes $1,3,4,4$; use $A_4/V_4\cong\Z/3\Z$.)
\end{exercise}

\begin{exercise}[$\star\star$]\label{ex:centralchar}
Let $\rho$ be irreducible of dimension $d$. Show that $\rho(z)$ is a scalar for every $z$ in the centre $Z(G)$, and deduce that $d^2\le[G:Z(G)]$. What does this give for the quaternion group $Q_8$, and what is its character table?
\end{exercise}

\begin{exercise}[$\star\star$]\label{ex:circulant}
For $f,g:\Z/n\Z\to\C$, let $(f*g)(k)=\sum_jf(j)g(k-j)$.
\begin{enumerate}[label=(\alph*)]
\item Show that $\widehat{f*g}(\chi)=n\,\hat f(\chi)\hat g(\chi)$.
\item A \emph{circulant} matrix has entries $C_{kj}=c(k-j)$ for some $c:\Z/n\Z\to\C$. Show that each character $\chi_m$, as a vector $(\chi_m(k))_k$, is an eigenvector of $C$, and find the eigenvalues.
\item A walker on $\Z/n\Z$ ($n$ odd) moves $\pm1$ with probability $\frac12$ each. Find the Fourier coefficients of its distribution after $t$ steps, starting from $0$, and show that it converges to the uniform distribution.
\end{enumerate}
\end{exercise}

\begin{exercise}[$\star\star$]\label{ex:dirichlet8}
Write the table of Dirichlet characters modulo $8$ and modulo $12$. Check Proposition~\ref{prop:dirichletorth} for $q=8$, $a=3$, and show that all characters modulo $8$ are real.
\end{exercise}

\begin{exercise}[$\star\star\star$]\label{ex:a5}
Find the character table of $A_5$ (Figure~\ref{fig:tables}), using its permutation representations on $\{1,\dots,5\}$ and on the $10$ two-element subsets, and the orthogonality relations.
\end{exercise}

\begin{exercise}[$\star\star\star$]\label{ex:divides}
A complex number is an \emph{algebraic integer} if it is a root of a monic polynomial with integer coefficients.
\begin{enumerate}[label=(\alph*)]
\item Show that $x\in\C$ is an algebraic integer iff $xM\subseteq M$ for some non-zero finitely generated subgroup $M$ of $(\C,+)$. Deduce that sums and products of algebraic integers are algebraic integers, and that a rational algebraic integer is an integer.
\item Let $\chi$ be irreducible of dimension $d$, $C$ a conjugacy class and $g\in C$. Show that $\sum_{x\in C}\rho(x)$ is the scalar $\omega_C=|C|\chi(g)/d$, and that $\omega_C$ is an algebraic integer.
\item Deduce that $d$ divides $|G|$.
\end{enumerate}
\end{exercise}

\begin{exercise}[$\star\star\star$]\label{pb:glnz}
(Finite groups of integer matrices.)
\begin{enumerate}[label=(\alph*)]
\item Show that $\cos\frac{2\pi}m$ is rational only for $m\in\{1,2,3,4,6\}$, and that every element of finite order of
$GL_2(\Z)$ has order $1$, $2$, $3$, $4$ or $6$.
\item Let $p\ge3$ be a prime and $A\in GL_n(\Z)$ with $A\equiv I_n\pmod p$ and $A\ne I_n$. Show that $A$ has infinite order.
Deduce that every finite subgroup of $GL_n(\Z)$ is isomorphic to a subgroup of $GL_n(\mathbb F_p)$, and that for $n=2$ its
order divides $48$.
\item Show that every finite subgroup of $GL_2(\Z)$ is cyclic or dihedral, of order at most $12$, and that $12$ is attained.
\end{enumerate}
\end{exercise}

\begin{exercise}[$\star\star$]\label{pb:commutators}
Let $G$ be a finite group and $N(g)$ the number of pairs $(x,y)\in G^2$ with $xyx^{-1}y^{-1}=g$.
\begin{enumerate}[label=(\alph*)]
\item (Frobenius, 1896.) Show that $N(g)=|G|\sum_\chi\frac{\chi(g)}{\chi(1)}$, the sum over the irreducible characters.
\item Deduce from the character table of $A_5$ (Exercise~11.7) that every element of $A_5$ is a commutator,
and compute $N(g)$ on each class.
\end{enumerate}
\end{exercise}

\begin{exercise}[$\star\star\star$]\label{pb:frobschur}
Let $G$ be finite, $r(g)$ the number of $x\in G$ with $x^2=g$, and for an irreducible character $\chi$ let
$\nu(\chi)=\frac1{|G|}\sum_{x\in G}\chi(x^2)$.
\begin{enumerate}[label=(\alph*)]
\item Show that $\nu(\chi)\in\{-1,0,1\}$, and that $\nu(\chi)\ne0$ if and only if $\chi$ is real-valued.
\item (Frobenius and Schur, 1906.) Show that $r=\sum_\chi\nu(\chi)\chi$.
\item Show that $\nu=-1$ for the irreducible character of degree $2$ of the quaternion group $Q_8$.
\item Show that $\nu(\chi)=1$ when the representation can be realised by real matrices. Every irreducible representation of
$\mathfrak S_n$ can be (Young, 1900). Deduce that the sum of the degrees of the irreducible characters of $\mathfrak S_n$ is
the number of $\sigma$ with $\sigma^2=e$.
\end{enumerate}
\end{exercise}

\section{Hints}

\begin{hint}{ex:pairs}
Count the fixed pairs of each class, then take inner products with the five rows of the table.
\end{hint}
\begin{hint}{ex:onedim}
$\C^*$ is abelian. For $\mathfrak S_n$, $A_4$ and $A_5$, find the derived subgroups with Chapters~4 and 5.
\end{hint}
\begin{hint}{ex:a4}
Three one-dimensional characters come from $\Z/3\Z$; then $12=1+1+1+9$. For the last row use column orthogonality, or the permutation representation on $4$ points.
\end{hint}
\begin{hint}{ex:centralchar}
Schur. Then $1=\langle\chi,\chi\rangle\ge\frac1{|G|}\sum_{z\in Z(G)}|\chi(z)|^2$. For $Q_8$: $|Z|=2$, $8=1+1+1+1+4$.
\end{hint}
\begin{hint}{ex:circulant}
(a) Expand and substitute $k=j+l$. (b) $(C\chi_m)(k)=\sum_jc(k-j)\chi_m(j)$; substitute $j=k-l$. (c) The distribution after $t$ steps is a $t$-fold convolution.
\end{hint}
\begin{hint}{ex:dirichlet8}
$(\Z/8\Z)^*\cong(\Z/2\Z)^2$, generated by $3$ and $5$; $(\Z/12\Z)^*=\{1,5,7,11\}$, likewise.
\end{hint}
\begin{hint}{ex:a5}
The first gives $\psi_4$, the second $\psi_4+\psi_5+\mathbf1$. Then $60=1+16+25+d^2+d'^2$; use column orthogonality on each column, and the fact that the values on $5$-cycles are real ($g$ and $g^{-1}$ are conjugate in $A_5$).
\end{hint}
\begin{hint}{ex:divides}
(a) If $M=\Z m_1+\dots+\Z m_k$ and $xm_i=\sum a_{ij}m_j$, then $x$ is an eigenvalue of the integer matrix $(a_{ij})$. (b) The products $e_Ce_D$, $e_C=\sum_{x\in C}\rho(x)$, are integer combinations of the $e_E$. (c) $\sum_C\omega_C\overline{\chi(g_C)}=|G|/d$.
\end{hint}

\begin{hint}{pb:glnz}
(a) $2\cos\frac{2\pi}m$ is an algebraic integer, and so is the trace of a matrix of finite order. (b) Reduce to $A$ of prime
order $q$, write $A=I+p^kB$ with $p\nmid B$, and expand $A^q$. (c) Average a scalar product over the group.
\end{hint}

\begin{hint}{pb:commutators}
Compute $\langle N,\chi\rangle$ with Schur's lemma: $\sum_y\rho(y)M\rho(y)^{-1}$ is a scalar.
\end{hint}

\begin{hint}{pb:frobschur}
(a) $\chi(x^2)$ is the difference of the characters of $\operatorname{Sym}^2V$ and $\Lambda^2V$. (b) Compute $\langle r,\chi\rangle$.
\end{hint}

\section{Solutions}

\begin{solution}{ex:pairs}
\idea A permutation representation has a particularly easy character: $\chi(g)$ is the number of basis vectors left fixed by $g$. Here the basis vectors are the six unordered pairs. Count fixed pairs in each conjugacy class, then take the inner products with the irreducible characters.

\writeup Let $V$ be the six-dimensional complex vector space with basis $e_{\{i,j\}}$ for $1\le i<j\le4$. The action of $\mathfrak S_4$ sends $e_{\{i,j\}}$ to $e_{\{g(i),g(j)\}}$. The trace counts the pairs left unchanged as \emph{sets}.

On the five conjugacy types $1,(12),(12)(34),(123),(1234)$, with class sizes $1,6,3,8,6$, the counts are:
\[
\begin{array}{c|ccccc}
\text{type}&1&(12)&(12)(34)&(123)&(1234)\\\hline
\text{fixed pairs}&6&2&2&0&0
\end{array}.
\]
For either involution, the two fixed pairs are $\{1,2\}$ and $\{3,4\}$. A $3$-cycle leaves no two-element subset unchanged; neither does a $4$-cycle. Thus
\[
\chi=(6,2,2,0,0).
\]

Use the irreducible character rows $\mathbf1$, $\varepsilon$, $\chi_2$, $\chi_3$, and $\chi'_3$ of the table in this chapter. Their multiplicities in $V$ are the inner products with $\chi$:
\[
\begin{aligned}
\langle\chi,\mathbf1\rangle&=\frac{6+6(2)+3(2)}{24}=1,\\
\langle\chi,\varepsilon\rangle&=\frac{6-6(2)+3(2)}{24}=0,\\
\langle\chi,\chi_2\rangle&=\frac{6(2)+3(2)(2)}{24}=1,\\
\langle\chi,\chi_3\rangle&=\frac{6(3)+6(2)(1)+3(2)(-1)}{24}=1,\\
\langle\chi,\chi_3'\rangle&=\frac{6(3)+6(2)(-1)+3(2)(-1)}{24}=0.
\end{aligned}
\]
Therefore
\[
V\cong\mathbf1\oplus V_2\oplus V_3.
\]
The dimensions provide a final check: $1+2+3=6$. The single trivial summand is the constant vector $\sum_{i<j}e_{\{i,j\}}$, which is visibly fixed by every permutation.
\end{solution}

\begin{solution}{ex:onedim}
\idea A one-dimensional representation is a multiplicative complex-valued homomorphism. Complex multiplication is commutative, so every commutator must act trivially. Therefore a one-dimensional representation only sees the abelianisation $G/G'$.

\writeup Let $\lambda:G\to GL_1(\C)=\C^*$ be a one-dimensional representation. For any $x,y\in G$,
\[
\lambda(xyx^{-1}y^{-1})
=\lambda(x)\lambda(y)\lambda(x)^{-1}\lambda(y)^{-1}=1.
\]
Thus every commutator is in $\ker\lambda$, and so is the subgroup $G'$ that they generate. The first isomorphism/correspondence properties of quotients show that $\lambda$ factors uniquely as
\[
G\longrightarrow G/G'\overset{\bar\lambda}{\longrightarrow}\C^*.
\]
Conversely, composing any homomorphism $\bar\lambda:G/G'\to\C^*$ with the quotient map gives a one-dimensional representation of $G$. These assignments are inverse.

The finite group $G/G'$ is abelian. By the structure theorem for finite abelian groups, it is a product of cyclic groups, and a cyclic factor of order $d$ has exactly $d$ complex characters, one for each $d$-th root of unity. Taking products shows that a finite abelian group $A$ has exactly $|A|$ one-dimensional characters. Hence the number sought for $G$ is $|G/G'|=[G:G']$.

\emph{Symmetric groups.} For $n\ge3$, the sign homomorphism shows $G'\subseteq A_n$. On the other hand, commutators of transpositions with a common point yield three-cycles. In fact, for distinct $a,b,c$,
\[
[(a\,b),(a\,c)]
=((a\,b)(a\,c))^2,
\]
which is a three-cycle or its inverse. Every three-cycle is obtained by choosing the labels appropriately. Because three-cycles generate $A_n$, we have $G'=A_n$. Thus $G/G'\cong\Z/2\Z$, so $\mathfrak S_n$ has exactly two one-dimensional representations: the trivial one and sign. For $n=2$ the group itself is cyclic of order $2$, giving the same answer.

\emph{Alternating group $A_4$.} Its normal subgroup
\[
V_4=\{e,(12)(34),(13)(24),(14)(23)\}
\]
has quotient $A_4/V_4\cong\Z/3\Z$, which is abelian; hence $A_4'\subseteq V_4$. The group $A_4$ is non-abelian, so its derived subgroup is not trivial. A non-trivial normal subgroup inside $V_4$ must contain one double transposition. All three double transpositions are conjugate in $A_4$, so normality forces it to contain all three. Thus $A_4'=V_4$, and there are $|A_4/V_4|=3$ one-dimensional representations.

\emph{Alternating group $A_5$.} Its derived subgroup $A_5'$ is normal. It is non-trivial because $A_5$ is not abelian. Since $A_5$ is simple, $A_5'=A_5$. The abelianisation is trivial, so the only one-dimensional representation is the trivial one.

The answers are therefore
\[
\boxed{\mathfrak S_n\ (n\ge2):2,\qquad A_4:3,\qquad A_5:1.}
\]
\end{solution}

\begin{solution}{ex:a4}
\idea Lift the characters of $\Z/3\Z$, then let orthogonality give the last row.
\writeup Let $\omega=e^{2i\pi/3}$. The classes are $e$, the double transpositions ($V_4\setminus\{e\}$, $3$ elements), the class of $(1\,2\,3)$ and that of $(1\,3\,2)$ ($4$ each). The quotient $A_4\to A_4/V_4\cong\Z/3\Z$ sends $(1\,2\,3)$ to a generator; the three characters of $\Z/3\Z$ give $\lambda_k=(1,1,\omega^k,\omega^{2k})$, $k=0,1,2$. The remaining irreducible has $d^2=12-3=9$, $d=3$. It is the permutation character $(4,0,1,1)$ minus the trivial one, $\chi=(3,-1,0,0)$, of norm $\frac{9+3+0+0}{12}=1$:
\begin{center}
\small
\begin{tabular}{c|cccc}
& $e$ & $(12)(34)$ & $(123)$ & $(132)$\\
size & $1$ & $3$ & $4$ & $4$\\\hline
$\lambda_0$ & $1$ & $1$ & $1$ & $1$\\
$\lambda_1$ & $1$ & $1$ & $\omega$ & $\omega^2$\\
$\lambda_2$ & $1$ & $1$ & $\omega^2$ & $\omega$\\
$\chi$ & $3$ & $-1$ & $0$ & $0$
\end{tabular}
\end{center}
The $3$-dimensional representation is the rotation group of the regular tetrahedron.
\end{solution}

\begin{solution}{ex:centralchar}
\idea A central element commutes with everything, so Schur makes it a scalar of modulus $1$.
\writeup For $z\in Z(G)$, $\rho(z)\rho(g)=\rho(zg)=\rho(gz)=\rho(g)\rho(z)$, so $\rho(z)$ is an intertwiner of the irreducible $V$, a scalar $\lambda I$, with $|\lambda|=1$ (a root of unity). So $|\chi(z)|=d$, and
\[
1=\langle\chi,\chi\rangle\ge\frac1{|G|}\sum_{z\in Z(G)}|\chi(z)|^2=\frac{|Z(G)|\,d^2}{|G|}.
\]
For $Q_8=\{\pm1,\pm i,\pm j,\pm k\}$: $Z=\{\pm1\}$, so $d^2\le4$. The five classes are
\[
\{1\},\{-1\},\{\pm i\},\{\pm j\},\{\pm k\}
\]
($i$ and $-i=jij^{-1}$ are conjugate); $8=\sum d_i^2$ with five terms $d_i\in\{1,2\}$ forces $1,1,1,1,4$. $Q_8/\{\pm1\}\cong(\Z/2\Z)^2$ gives four one-dimensional characters: $\pm1$ on the classes of $i,j,k$ with an even number of minus signs. The last one, $\chi$, has $\chi(1)=2$. Its other values follow from column orthogonality with the column of $1$ (Corollary~\ref{cor:columns}): the four one-dimensional characters are $1$ at $-1$, so $4+2\chi(-1)=0$ and $\chi(-1)=-2$; and at $i$ they sum to $1+1-1-1=0$, so $\chi(i)=0$, and likewise $\chi(j)=\chi(k)=0$. It is realised by the quaternions of Chapter~12 acting on $\C^2$. Curiously, $D_4$ has the same character table, although $D_4\not\cong Q_8$: the character table does not determine the group.
\end{solution}

\begin{solution}{ex:circulant}
\idea Convolution is the algebraic rule for adding independent positions on a cyclic group. Fourier characters turn convolution into scalar multiplication, which makes repeated steps easy to compute. For the convergence rate, compare \emph{absolute values} of the Fourier eigenvalues: the most negative one can decay more slowly than the first positive one.

\writeup Write $\chi_m(k)=e^{2\pi i mk/n}$ for $m,k\in\Z/n\Z$, and use the normalised Fourier transform
\[
\hat f(\chi)=\frac1n\sum_{k\in\Z/n\Z}f(k)\overline{\chi(k)}.
\]

\emph{(a) Transform of a convolution.} Starting from the definitions and substituting $\ell=k-j$, we get
\[
\begin{aligned}
n\,\widehat{f*g}(\chi)
&=\sum_k\sum_j f(j)g(k-j)\overline{\chi(k)}\\
&=\sum_j\sum_\ell f(j)g(\ell)\overline{\chi(j+\ell)}\\
&=\left(\sum_j f(j)\overline{\chi(j)}\right)
  \left(\sum_\ell g(\ell)\overline{\chi(\ell)}\right)\\
&=n^2\hat f(\chi)\hat g(\chi).
\end{aligned}
\]
We used the character law $\chi(j+\ell)=\chi(j)\chi(\ell)$ in the third line. Dividing by $n$ proves
\[
\boxed{\widehat{f*g}(\chi)=n\,\hat f(\chi)\hat g(\chi)}.
\]

\emph{(b) Circulant eigenvectors.} Let $C_{kj}=c(k-j)$. Then
\[
\begin{aligned}
(C\chi_m)(k)
&=\sum_jc(k-j)\chi_m(j)\\
&=\sum_\ell c(\ell)\chi_m(k-\ell)\\
&=\chi_m(k)\sum_\ell c(\ell)\overline{\chi_m(\ell)}.
\end{aligned}
\]
Thus $\chi_m$ is an eigenvector with eigenvalue
\[
\boxed{\lambda_m=n\hat c(\chi_m)
=\sum_{\ell=0}^{n-1}c(\ell)e^{-2\pi i m\ell/n}}.
\]
The $n$ characters form an orthogonal basis, so this lists all the eigenvalues, with multiplicity.

\emph{(c) The random walk.} A single step has law
\[
s=\tfrac12(\delta_1+\delta_{-1}).
\]
Starting from $0$, after $t$ independent steps the position law is the $t$-fold convolution $s^{*t}$. The transform of one step is
\[
n\hat s(\chi_m)=\tfrac12(e^{-2\pi im/n}+e^{2\pi im/n})
=\cos\frac{2\pi m}{n}.
\]
Applying (a) repeatedly, including $t=0$, gives
\[
\boxed{\widehat{s^{*t}}(\chi_m)
=\frac1n\left(\cos\frac{2\pi m}{n}\right)^t}.
\]
Fourier inversion consequently yields the exact probability
\[
s^{*t}(k)=\frac1n\sum_{m=0}^{n-1}
\left(\cos\frac{2\pi m}{n}\right)^t
e^{2\pi i mk/n}.
\]

The mode $m=0$ equals $1/n$, the uniform distribution. For odd $n\ge3$, none of the other angles $2\pi m/n$ is a multiple of $\pi$, hence
\[
\left|\cos\frac{2\pi m}{n}\right|<1
\quad (m\ne0).
\]
Every non-constant mode therefore tends to zero, proving $s^{*t}(k)\to1/n$ for every $k$.

\emph{Which mode is slowest?} The negative eigenvalues matter as much as the positive ones. For odd $n$,
\[
\max_{1\le m<n}\left|\cos\frac{2\pi m}{n}\right|
=\cos\frac{\pi}{n}.
\]
The maximum occurs at $m=(n-1)/2$ and $m=(n+1)/2$, whose eigenvalues are $-\cos(\pi/n)$. Thus the slowest non-constant contribution is of magnitude
\[
\cos^t\frac{\pi}{n}
=\exp\!\left(-\frac{\pi^2t}{2n^2}
+O\!\left(\frac{t}{n^4}\right)\right)
\]
for large odd $n$. The negative sign makes that mode alternate with $t$, but its magnitude decays. This is why the natural relaxation scale is of order $n^2$.

For even $n$, the character $m=n/2$ has eigenvalue $\cos\pi=-1$. That component never decays: the walk alternates between the two parity classes, so it does not converge to the uniform law without an aperiodic modification such as a holding probability.
\end{solution>

\begin{solution}{ex:dirichlet8}
\idea A Dirichlet character modulo $q$ is a character of the unit group $(\Z/q\Z)^*$, extended by $0$ on non-units. For $q=8$ and $12$, every non-identity unit has order $2$. The character is therefore determined by assigning signs to two independent generators.

\writeup Modulo $8$, the units are
\[
(\Z/8\Z)^*=\{1,3,5,7\}.
\]
They satisfy $3^2=5^2=7^2=1$ and $3\cdot5=7$ modulo $8$. Thus the group is $(\Z/2\Z)^2$, with independent generators $3$ and $5$. For a character $\chi$, multiplicativity forces
\[
\chi(3)^2=\chi(5)^2=1,
\quad\text{so}\quad \chi(3),\chi(5)\in\{\pm1\},
\qquad\chi(7)=\chi(3)\chi(5).
\]
Each independent pair of signs gives exactly one character. Here is the full table on the units:
\begin{center}
\small
\begin{tabular}{c|rrrr}
modulo $8$ & $1$ & $3$ & $5$ & $7$\\\hline
$\chi_0$ & $1$ & $1$ & $1$ & $1$\\
$\chi_1$ & $1$ & $-1$ & $1$ & $-1$\\
$\chi_2$ & $1$ & $1$ & $-1$ & $-1$\\
$\chi_3$ & $1$ & $-1$ & $-1$ & $1$
\end{tabular}
\end{center}
On residues $0,2,4,6$, each of these Dirichlet characters has value $0$.

Modulo $12$, the units are $\{1,5,7,11\}$, with $5^2=7^2=1$ and $5\cdot7=11$. We obtain the same abstract group and four characters:
\begin{center}
\small
\begin{tabular}{c|rrrr}
modulo $12$ & $1$ & $5$ & $7$ & $11$\\\hline
$\psi_0$ & $1$ & $1$ & $1$ & $1$\\
$\psi_1$ & $1$ & $-1$ & $1$ & $-1$\\
$\psi_2$ & $1$ & $1$ & $-1$ & $-1$\\
$\psi_3$ & $1$ & $-1$ & $-1$ & $1$
\end{tabular}
\end{center}
The characters vanish at every residue not coprime to $12$.

We can now check the orthogonality projector of Proposition~\ref{prop:dirichletorth} for $q=8$ and $a=3$. All values are real, so complex conjugation has no effect. Evaluate
\[
\frac1{\varphi(8)}\sum_{\chi\bmod8}\overline{\chi(3)}\chi(n)
=\frac14\sum_{\chi\bmod8}\chi(3)\chi(n).
\]
The four products $\chi(3)\chi(n)$ sum to $0$ at $n=1,5,7$ and to $4$ at $n=3$; at even $n$ every character is zero. For instance at $n=3$ every product equals $\chi(3)^2=1$, while at $n=5$ the four products are $1,-1,-1,1$.

The projector is therefore $1$ exactly when $n\equiv3\pmod8$, and $0$ otherwise. Finally, the tables prove that every character modulo $8$ is real: all its values belong to $\{0,\pm1\}$. This is a consequence of the exponent-$2$ unit group, not a general property of Dirichlet characters.
\end{solution}

\begin{solution}{ex:a5}
\idea Two permutation representations give $\psi_4$ and $\psi_5$; orthogonality pins down the two missing rows, and the golden ratio appears as a root of $x^2-x-1$.
\writeup Classes (Theorem~\ref{thm:a5}): $e$, $(12)(34)$, $(123)$, $(12345)$, $(13524)$, of sizes $1,15,20,12,12$. On $\{1,\dots,5\}$ the fixed-point counts are $(5,1,2,0,0)$; minus $\mathbf1$, $\psi_4=(4,0,1,-1,-1)$, of norm $\frac{16+0+20+12+12}{60}=1$. On the $10$ pairs, fixed pairs: $10$; $2$ ($\{1,2\}$, $\{3,4\}$); $1$ ($\{4,5\}$); $0$; $0$. Minus $\mathbf1$ and $\psi_4$: $\psi_5=(5,1,-1,0,0)$, of norm $\frac{25+15+20}{60}=1$. The remaining irreducibles satisfy $\sum d^2=60-1-16-25=18$, with five classes, so there are two, of dimension $3$ ($9+9$): $\psi_3=(3,a,b,c,c')$ and $\psi_3'=(3,a',b',e,e')$.

Column orthogonality (Corollary~\ref{cor:columns}) between the column of $g$ and the column of $e$ gives $\sum_\chi\chi(e)\chi(g)=0$, and the column of $g$ has squared norm $|C_G(g)|=60/|\text{class}|$. The values are real: each $g\in A_5$ is conjugate to $g^{-1}$ (for the $5$-cycle, $(2\,5)(3\,4)$ conjugates $(12345)$ to $(15432)$), so $\chi(g)=\chi(g^{-1})=\overline{\chi(g)}$.
\begin{itemize}
\item $(12)(34)$, centraliser of order $4$: $1+0+5+3a+3a'=0$ and $1+0+1+a^2+a'^2=4$; so $a+a'=-2$, $a^2+a'^2=2$, $a=a'=-1$.
\item $(123)$, centraliser of order $3$: $1+4-5+3b+3b'=0$ and $1+1+1+b^2+b'^2=3$; so $b=b'=0$.
\item $(12345)$, centraliser of order $5$: $1-4+0+3c+3e=0$ and $1+1+0+c^2+e^2=5$; so $c+e=1$, $c^2+e^2=3$, $ce=-1$: $c,e$ are the roots of $x^2-x-1$, i.e.\ $\varphi$ and $1-\varphi$. The same holds for $c',e'$; and $\langle\psi_3,\mathbf1\rangle=0$ gives $3-15+0+12c+12c'=0$, $c+c'=1$, so $c'=1-c$.
\end{itemize}
Choosing $\psi_3(12345)=\varphi$: $\psi_3=(3,-1,0,\varphi,1-\varphi)$, $\psi_3'=(3,-1,0,1-\varphi,\varphi)$. $\psi_3$ is the rotation representation of $A_5$ as the symmetry group of the icosahedron (Chapter~12): a rotation by $\frac{2\pi}5$ has trace $1+2\cos\frac{2\pi}5=\varphi$.
\end{solution}

\begin{solution}{ex:divides}
\idea Integrality is stability of a finitely generated lattice; the central elements $e_C$ span such a lattice, and they act by scalars.
\writeup (a) If $x^n+a_{n-1}x^{n-1}+\dots+a_0=0$ with $a_i\in\Z$, then $M=\Z+\Z x+\dots+\Z x^{n-1}$ satisfies $xM\subseteq M$ and $1\in M$. Conversely, if $M=\Z m_1+\dots+\Z m_k\ne0$ and $xm_i=\sum_ja_{ij}m_j$ with $a_{ij}\in\Z$, then $(xI-A)m=0$ for the non-zero vector $m=(m_i)$, so $\det(xI-A)=0$: $x$ is a root of the monic integer polynomial $\det(XI-A)$. If $x,y$ are algebraic integers, of degrees $n,n'$ as above, $M=\sum_{i<n,j<n'}\Z x^iy^j$ is finitely generated, contains $1$, and is stable under multiplication by $x$, $y$, hence by $x+y$ and $xy$. A rational root $a/b$ (in lowest terms) of a monic integer polynomial has $b\mid1$ (Volume~1, rational roots).

(b) $e_C=\sum_{x\in C}\rho(x)$ commutes with every $\rho(g)$ (conjugation by $g$ permutes $C$), so it is a scalar $\omega_C$ (Schur); taking traces, $d\,\omega_C=|C|\chi(g)$. Summing over pairs,
\[
e_Ce_D=\sum_{x\in C,\,y\in D}\rho(xy)=\sum_Ea_{CDE}\,e_E,
\]
where $a_{CDE}\ge0$ is the number of pairs $(x,y)\in C\times D$ with $xy=z$, for a fixed $z\in E$ (the same for all $z\in E$, by conjugation). So $\omega_C\omega_D=\sum_Ea_{CDE}\,\omega_E$. So $M=\sum_E\Z\omega_E$ is finitely generated, non-zero ($\omega_{\{e\}}=1$), and $\omega_CM\subseteq M$: $\omega_C$ is an algebraic integer by (a).

(c) $\overline{\chi(g)}$ is a sum of roots of unity, an algebraic integer. Grouping by classes,
\[
\begin{aligned}\frac{|G|}d=\frac{|G|}d\langle\chi,\chi\rangle&=\frac1d\sum_C|C|\chi(g_C)\overline{\chi(g_C)}\eqbr=\sum_C\omega_C\overline{\chi(g_C)},\end{aligned}
\]
an algebraic integer by (a), and rational. So $|G|/d\in\Z$.
\end{solution}

\begin{solution}{pb:glnz}
\idea An integer matrix of finite order has eigenvalues that are roots of unity, and its trace is an integer: that leaves only
a few angles. Reducing modulo $p$ is a homomorphism, and the question is whether it can kill an element of finite order; the
$p$-adic size of $A-I$ shows it cannot, because raising $I+p^kB$ to a power $q$ multiplies the ``error'' $p^kB$ by $q$ to
first order and only adds higher powers of $p$. Finally, a finite group of real matrices preserves some scalar product,
so it is a group of rotations and reflections in disguise.

\writeup (a) $2\cos\frac{2\pi}m=\zeta+\zeta^{-1}$ with $\zeta=e^{2i\pi/m}$, a sum of algebraic integers, hence an algebraic
integer (Volume~3, Chapter~11). If it is rational, it is an integer (Volume~2, Chapter~1), in $[-2,2]$, so $\cos\frac{2\pi}m\in
\{0,\pm\frac12,\pm1\}$ and $m\in\{1,2,3,4,6\}$.

Let $A\in GL_2(\Z)$ with $A^m=I$. Its minimal polynomial divides $X^m-1$, which has distinct roots in $\C$, so $A$ is
diagonalisable over $\C$ (Volume~1, Chapter~8), with eigenvalues $\lambda_1,\lambda_2$ that are roots of unity, and the order of
$A$ is the lcm of their orders. The characteristic polynomial has integer coefficients. If $\lambda_1$ is real, so is
$\lambda_2$, both are $\pm1$, and the order is $1$ or $2$. Otherwise $\lambda_2=\overline{\lambda_1}$, and
$\operatorname{tr}A=2\operatorname{Re}\lambda_1=2\cos\frac{2\pi a}{d}$ with $d$ the order of $\lambda_1$ and $\gcd(a,d)=1$.
This is an integer in $]-2,2[$, so the cosine is $0$ or $\pm\frac12$, and $d\in\{3,4,6\}$ (the angle $\frac{2\pi a}d$ is
$\pm\frac\pi2$, $\pm\frac\pi3$ or $\pm\frac{2\pi}3$ modulo $2\pi$). So the order of $A$ is in $\{1,2,3,4,6\}$.

(b) Suppose $A$ has finite order $m>1$, and let $q$ be a prime dividing $m$; then $A'=A^{m/q}$ has order $q$ and still
$A'\equiv I\pmod p$. Rename it $A$. Write $A=I+p^kB$ with $k\ge1$ and $B\in M_n(\Z)$ not divisible by $p$. Then
\[
I=A^q=I+qp^kB+\sum_{j=2}^q\binom qjp^{jk}B^j,\qquad\text{so}\qquad qp^kB=-\sum_{j=2}^q\binom qjp^{jk}B^j.
\]
If $q\ne p$: the right-hand side is divisible by $p^{2k}$, while $qp^kB$ is divisible by $p^k$ but not by $p^{k+1}$ (some
entry of $B$ is prime to $p$, and so is $q$). As $2k\ge k+1$, this is impossible. If $q=p$: for $2\le j<p$, $p\mid\binom pj$, so
the term is divisible by $p^{1+jk}$, with $1+jk\ge k+2$; the last term is divisible by $p^{pk}$, and $pk\ge k+2$ because
$p\ge3$. So the right-hand side is divisible by $p^{k+2}$, and $p^{k+1}B$ is not: impossible again.

So the reduction map $GL_n(\Z)\to GL_n(\mathbb F_p)$, a group homomorphism, is injective on every finite subgroup $G$, whose
order therefore divides $|GL_n(\mathbb F_p)|$ (Lagrange). For $n=2$, $|GL_2(\mathbb F_3)|=8\cdot6=48$, so $|G|$ divides $48$.
(For every $p\ge5$, $|GL_2(\mathbb F_p)|=p(p-1)^2(p+1)$ is a multiple of $48$, so other primes add nothing.)

(c) Let $G\subset GL_2(\Z)$ be finite. For $x,y\in\R^2$ put $\langle x,y\rangle_G=\sum_{g\in G}(gx)\cdot(gy)$. This is a scalar
product (positive definite, as a sum of positive definite forms), and it is $G$-invariant: $\langle hx,hy\rangle_G=\langle x,
y\rangle_G$ for $h\in G$, by reindexing $g\mapsto gh$. In an orthonormal basis for $\langle\,,\rangle_G$ the elements of $G$
become orthogonal matrices, so $G$ is isomorphic to a finite subgroup of $O(2)$. Let $G^+$ be its subgroup of rotations
(determinant $1$), of index $1$ or $2$. Rotations form a group isomorphic to the unit circle of $\C$, whose finite subgroups are
cyclic (Volume~3, Chapter~1), so $G^+$ is cyclic of some order $m$, generated by a rotation of order $m$, and by (a)
$m\in\{1,2,3,4,6\}$. If $G\ne G^+$, let $s\in G$ be a reflection; then $G=G^+\cup sG^+$ and $srs^{-1}=r^{-1}$ for every rotation
$r$, so $G$ is dihedral of order $2m$. Hence $|G|\le12$. The order $12$ is attained: $A=\begin{psmallmatrix}0&-1\\1&1
\end{psmallmatrix}$ (characteristic polynomial $X^2-X+1$, order $6$) and $B=\begin{psmallmatrix}0&1\\1&0\end{psmallmatrix}$
satisfy $BAB^{-1}=A^{-1}$ and generate a dihedral group of order $12$, the symmetries of the hexagonal lattice.
\tools Volume~3, Chapters~1 and~11 (finite subgroups of $\C^*$, algebraic integers); Volume~2, Chapters~1 and~3 (rational
algebraic integers, congruences, binomial coefficients modulo $p$); Volume~1, Chapter~8 (diagonalisation, scalar products).
\leads This is the \emph{crystallographic restriction}: a lattice in the plane has no symmetry of order $5$ or $7$, which is
why quasicrystals (Shechtman, Nobel prize 2011) were a shock. Minkowski (1887) used (b) for all primes at once to find the
largest finite subgroups of $GL_n(\Q)$, and Bieberbach's theorems (1911) classify crystallographic groups in all dimensions.
\end{solution}

\begin{solution}{pb:commutators}
\idea $N$ is a class function, so it is determined by its inner products with the irreducible characters. Summing $\chi$ over all
commutators $[x,y]$, the sum over $y$ of $\rho(y)\rho(x)^{-1}\rho(y)^{-1}$ averages $\rho(x)^{-1}$ into a scalar, by Schur's lemma.

\writeup (a) If $[x,y]=g$ then $[hxh^{-1},hyh^{-1}]=hgh^{-1}$, so $N$ is a class function. Let $\rho$ be irreducible with
character $\chi$ of degree $d$. For any matrix $M$, $\sum_y\rho(y)M\rho(y)^{-1}$ commutes with every $\rho(z)$, so it is a scalar
(Schur's lemma, Volume~3, Chapter~11), and comparing traces it is $\frac{|G|\operatorname{tr}M}dI$. With $M=\rho(x)^{-1}$:
\[
\sum_y\chi(xyx^{-1}y^{-1})=\operatorname{tr}\Bigl(\rho(x)\sum_y\rho(y)\rho(x)^{-1}\rho(y)^{-1}\Bigr)=\frac{|G|}d\chi(x)\chi(x^{-1})=\frac{|G|}d|\chi(x)|^2.
\]
Summing over $x$, $\sum_{x,y}\chi([x,y])=\frac{|G|^2}d$, which is real; so
$\langle N,\chi\rangle=\frac1{|G|}\sum_gN(g)\overline{\chi(g)}=\frac1{|G|}\overline{\sum_{x,y}\chi([x,y])}=\frac{|G|}d$. The irreducible characters
are an orthonormal basis of class functions (Theorem~11.10), so $N=\sum_\chi\frac{|G|}{\chi(1)}\chi$.

(b) The table of $A_5$ on the classes $e$, $(12)(34)$, $(123)$, $(12345)$, $(13524)$ has rows $(1,1,1,1,1)$,
$(3,-1,0,\phi,1-\phi)$, $(3,-1,0,1-\phi,\phi)$, $(4,0,1,-1,-1)$, $(5,1,-1,0,0)$, with $\phi=\frac{1+\sqrt5}2$. So
\[
\frac{N(g)}{60}=1+\frac{\chi_2(g)+\chi_3(g)}3+\frac{\chi_4(g)}4+\frac{\chi_5(g)}5
\]
is $5$ at $e$, $1-\frac23+\frac15=\frac8{15}$ at $(12)(34)$, $1+\frac14-\frac15=\frac{21}{20}$ at $(123)$, and $1+\frac13-\frac14=\frac{13}{12}$ at
both $5$-cycles (since $\phi+1-\phi=1$). So $N=300,32,63,65,65$: never $0$. Check: $300+15\cdot32+20\cdot63+24\cdot65=3600=60^2$.
\tools Volume~3, Chapter~11 (Schur's lemma, orthogonality, the table of $A_5$).
\leads $N(e)=|G|\cdot k(G)$ recovers the count of commuting pairs. Ore conjectured in 1951 that every element of every finite
non-abelian simple group is a commutator; Liebeck, O'Brien, Shalev and Tiep proved it in 2010, with this formula as the main tool:
they show that $\sum_\chi\frac{\chi(g)}{\chi(1)}\ne0$ by bounding character values.
\end{solution}

\begin{solution}{pb:frobschur}
\idea $\chi(x^2)$ is the trace of $\rho(x)^2$, which is the difference of the characters of the symmetric and alternating squares.
So $\nu(\chi)$ is the difference of the dimensions of their invariants, and $V\otimes V$ has at most one invariant line. For (b),
$r$ is a class function and its inner products with characters are the $\nu$'s.

\writeup (a) If $\rho(x)$ has eigenvalues $\lambda_1,\dots,\lambda_d$, then $\operatorname{Sym}^2V$ has eigenvalues $\lambda_i\lambda_j$,
$i\le j$, and $\Lambda^2V$ has $\lambda_i\lambda_j$, $i<j$. So their characters are $\frac{\chi(x)^2\pm\chi(x^2)}2$, and
$\chi(x^2)=\chi_S(x)-\chi_\Lambda(x)$. Hence $\nu(\chi)=\langle\chi_S,1\rangle-\langle\chi_\Lambda,1\rangle=s-a$, where $s,a\ge0$ are the
dimensions of the invariant vectors. Their sum is $\langle\chi^2,1\rangle=\frac1{|G|}\sum\chi(x)^2=\langle\chi,\bar\chi\rangle$, which is
$1$ if $\bar\chi=\chi$ and $0$ otherwise ($\bar\chi$ is also an irreducible character, of the dual representation). So
$(s,a)\in\{(0,0),(1,0),(0,1)\}$ and $\nu\in\{0,1,-1\}$, with $\nu\ne0$ exactly when $\chi$ is real-valued.

(b) $r$ is a class function, and $\langle r,\chi\rangle=\frac1{|G|}\sum_gr(g)\overline{\chi(g)}=\frac1{|G|}\sum_x\overline{\chi(x^2)}=\nu(\chi)$,
real by (a). So $r=\sum\nu(\chi)\chi$.

(c) In $Q_8$, $x^2=1$ only for $x=\pm1$, so $r(1)=2$. The characters are four of degree $1$ (real, since $Q_8/\{\pm1\}\cong(\Z/2\Z)^2$, with
$\nu=1$) and one of degree $2$. So $2=4\cdot1+2\nu$, and $\nu=-1$.

(d) If $\rho$ is real, averaging any positive definite symmetric bilinear form over $G$ gives an invariant one, a non-zero vector
of $\operatorname{Sym}^2V^*\cong\operatorname{Sym}^2V$ (as $\chi$ is real, $V\cong V^*$). So $s=1$ and $\nu=1$. For $\mathfrak S_n$ all
$\nu$ are then $1$, and evaluating (b) at $e$: $\#\{\sigma:\sigma^2=e\}=r(e)=\sum_\chi\chi(1)$.
\tools Volume~3, Chapter~11 (characters, orthogonality); Volume~1, Chapter~8 (bilinear forms, eigenvalues).
\leads For $n=4$: degrees $1,1,2,3,3$, sum $10$, and $\mathfrak S_4$ has $1+6+3=10$ involutions with the identity. Since the degrees count
standard Young tableaux, (d) says that there are as many standard Young tableaux with $n$ boxes as involutions; the
Robinson--Schensted correspondence gives a bijection. The value $\nu=-1$ marks the \emph{quaternionic} representations.
\end{solution}

\section*{Sources and further reading}

\begin{itemize}
\item G.~Frobenius, \emph{Über Gruppencharaktere}, Sitzungsber.\ Preuss.\ Akad.\ Wiss.\ Berlin (1896), 985--1021. The birth of character theory.
\item H.~Maschke, \emph{Beweis des Satzes, dass diejenigen endlichen linearen Substitutionsgruppen, in welchen einige durchgehends verschwindende Coefficienten auftreten, intransitiv sind}, Math.\ Ann.\ 52 (1899), 363--368. Complete reducibility.
\item I.~Schur, \emph{Neue Begründung der Theorie der Gruppencharaktere}, Sitzungsber.\ Preuss.\ Akad.\ Wiss.\ Berlin (1905), 406--432. Schur's lemma and the proofs of this chapter.
\item J.-P.~Serre, \emph{Linear Representations of Finite Groups}, Springer, 1977, Part~I. The model for this chapter, in forty pages.
\item G.~James and M.~Liebeck, \emph{Representations and Characters of Groups}, 2nd ed., Cambridge University Press, 2001. Many character tables worked out, including $A_5$.
\item P.~Diaconis, \emph{Group Representations in Probability and Statistics}, IMS Lecture Notes 11, 1988. Characters and Fourier analysis for random walks and shuffles.
\item J.~W.~Cooley and J.~W.~Tukey, \emph{An algorithm for the machine calculation of complex Fourier series}, Math.\ Comp.\ 19 (1965), 297--301. The fast Fourier transform.
\item P.~G.~L.~Dirichlet, \emph{Beweis des Satzes, dass jede unbegrenzte arithmetische Progression \dots\ unendlich viele Primzahlen enthält}, Abh.\ Königl.\ Preuss.\ Akad.\ Wiss.\ (1837), 45--81. Characters of $(\Z/q\Z)^*$, before the word existed.
\end{itemize}
